\assignq{4}{2}{Zeeman slower: scattering limit, resonance profile and panels (b)--(e)}{1--2}
Consider a beam propagating along \(+z\), illuminated by
a counter-propagating laser. Let \(\Gamma=1/\tau\)
be the excited-state \emph{population decay rate},
\(s=I/I_{\rm sat}\) the saturation parameter and
\(\Delta_{\rm eff}\) the laser detuning in the
atom's rest frame.

\subsection*{(a) Maximum scattering rate and design acceleration}
The two-level scattering rate is
\[
R_{\rm sc}=\frac{\Gamma}{2}
\frac{s}{1+s+(2\Delta_{\rm eff}/\Gamma)^2}.
\]
Near saturation the maximum is
\(R_{\max}=\Gamma/2\).
Every photon absorbed from the counter-propagating
laser changes the atom's directed momentum by
approximately \(-\hbar k\). Averaging over random
spontaneous recoil gives
\[
\boxed{F_{\max}=\hbar k\Gamma/2,\quad
a_{\max}=\frac{\hbar k\Gamma}{2M}}.
\]
A robust real slower typically designs for
\(a=\eta a_{\max}\), where \(0<\eta<1\),
leaving a margin for power fluctuations, finite
linewidth and deviations from exact resonance.
Thus, for incident speed \(v_0\), the required
ideal length is \(L_0=v_0^2/(2a)\).

\subsection*{(b) Why the optimal field is square-root shaped}
Choose a Zeeman shift
\(\mu' B(z)/\hbar\), where \(\mu'\) is the
effective differential magnetic moment,
such that the changing resonance frequency
compensates the Doppler shift. With a suitable
sign convention, the resonance condition is
\[
\Delta_L+k v(z)-\frac{\mu' B(z)}{\hbar}=0.
\]
Therefore \(B(z)=B_{\rm offset}
+(\hbar k/\mu')v(z)\) for fixed laser detuning.
At constant deceleration,
\[
v(z)^2=v_0^2-2az,\qquad
\boxed{v(z)=v_0\sqrt{1-z/L_0},\quad
B(z)=B_{\rm offset}+B_0\sqrt{1-z/L_0}},
\]
where \(B_0=\hbar k v_0/\mu'\) with the
appropriate sign. The square-root profile
keeps the photon force close to the
chosen constant \(Ma\) while the atom slows.

\subsection*{(c) What sets the capture velocity?}
Three main limits matter: the \emph{available
slower length}, the \emph{maximum sustainable
scattering force} and the \emph{magnetic-field
range}/laser detuning needed to bring a fast
atom into resonance. In the ideal fixed-acceleration
model,
\[
\boxed{v_{\rm capture}\lesssim
\sqrt{v_f^2+2aL}},
\]
subject to the entrance-resonance condition
\[
|\Delta_L+kv_0-\mu' B(0)/\hbar|
\lesssim\text{the power-broadened linewidth}.
\]
A longer slower cannot compensate an insufficient
field range or an atom that never enters the
laser-coupling resonance window.

\subsection*{(d) Analyse the supplied alternative field plots}
The original Assignment 4 page 2 draws
\(B(z)\) for a reference optimal field (a)
and four alternatives (b)--(e).
To analyse any alternative, first translate
\(B(z)\) into the \emph{resonant velocity}
\(v_{\rm res}(z)\propto B(z)-B_{\rm offset}\).
Its gradient determines the acceleration
needed to stay locked:
\[
\boxed{a_{\rm req}(z)=
-v_{\rm res}(z)\frac{\diff v_{\rm res}}{\diff z}.}
\]
The atom cannot follow a resonant trajectory
requiring \(a_{\rm req}>a_{\max}\), and a
\emph{positive} resonant-velocity slope requires
the atom to accelerate in the beam direction,
opposite to the counter-propagating laser force.

\begin{center}\small
\begin{tabularx}{\linewidth}{@{}p{8mm}p{34mm}X@{}}
\toprule
Panel&Observed field shape&Outcome and physical reason\\\midrule
(a)&Reference square root; zero at \(L_0\)&
Yes in the ideal model: constant allowed
deceleration \(a\), reaching \(v=0\) at \(L_0\).\\
(b)&Shallow straight line;
zero near \(2L_0\)&
Can slow to \emph{arbitrarily small} velocity
over a longer length. Taking
\(v_{\rm res}=v_0(1-z/(2L_0))\) gives
\(a_{\rm req}=a(1-z/(2L_0))\le a\).
The mathematical \(v=0\) endpoint is
approached asymptotically in time.\\
(c)&Steep straight line;
zero near \(L_0\)&
At its entrance the required deceleration is
\(2a\), twice that of reference (a).
If (a) was designed at the available force,
the atoms fall out of resonance and do not
follow the curve to rest.\\
(d)&A strong down--up--down excursion&
The rising-field section corresponds to
an \emph{increasing} resonant speed and cannot
be tracked by atoms under a purely opposing
scattering force. The first steep section can
also exceed the force budget. Slowing may
resume if the atom is later recaptured,
but full stopping is not guaranteed.\\
(e)&Rapid oscillations before the optimal tail&
The ideal resonance velocity oscillates rapidly,
demanding alternating acceleration/deceleration.
If oscillations exceed the effective linewidth
or force response, atoms are detuned and
lost; very small slow ripples may be tolerated.
Full slowing cannot be assured from the
nominal endpoint alone.\\
\bottomrule
\end{tabularx}
\end{center}

A useful additional check for (b): as
\(v(t)\propto e^{-ct}\), the atom has
\(v\to0\) only for \(t\to\infty\) within the
ideal linearly decreasing resonant-velocity
model. Real atoms undergo stochastic photon
recoil and exit at a nonzero capture velocity;
this makes the distinction between a
\emph{mathematical endpoint} and practical
capture physically relevant.

\methodnote{Always determine whether a
field profile lies within the available
\emph{force and linewidth} constraints.
A plot that reaches \(B=0\) does not on
its own guarantee that the atoms follow it
or end at rest.}
